\section{Related Work}
\label{sec:related-work}

OutlierMigrate sits at the intersection of post-training quantization,
rotation-based outlier suppression, long-decode reasoning quantization, and
state-space-model deployment.  The section is written to be independent of
the current Nemotron gate outcome: if later cross-model validation is a pass,
ambiguous, or a kill, the literature boundary remains the same.  The claim is
not that existing quantizers fail broadly; it is that decode-position
channel-set drift is a specific stressor for static protected-channel
assumptions, and any method claim must be conditioned on the audited gate that
supports it.

\paragraph{Post-training quantization and static outlier protection.}
Classical LLM PTQ methods establish the default deployment baseline.  GPTQ
casts weight-only quantization as a layer-wise reconstruction problem with
second-order error compensation \citep{frantar2022gptq}.  SmoothQuant moves
activation quantization difficulty into weights using calibration-time
activation statistics, making W8A8 inference practical for large models
\citep{xiao2023smoothquant}.\footnote{We use ``outlier migration'' only for
decode-position channel-set drift: the identity of high-magnitude activation
channels changing across generated positions in the same trace.  This differs
from SmoothQuant's activation-to-weight difficulty transfer and from
token- or bitwidth-conditioned uses of the term in other quantization papers.}
AWQ protects activation-salient weight channels and is a strong weight-only
baseline for W4A16-style deployment \citep{lin2023awq}.  KV-cache quantizers
attack a different memory object: KIVI uses asymmetric low-bit KV-cache
quantization with per-channel keys and per-token values \citep{liu2024kivi},
while KVQuant combines per-channel key quantization, pre-RoPE key handling,
non-uniform datatypes, and dense/sparse outlier handling for long-context KV
inference \citep{hooper2024kvquant}.  These systems motivate our controls, but
they do not by themselves answer whether a fixed activation-channel protection
map remains valid thousands of generated tokens after calibration.

\paragraph{Rotation and incoherence methods.}
Rotation-based quantizers reduce outliers by changing basis before low-bit
rounding.  QuaRot applies Hadamard rotations to make LLM inference more
outlier-free under low-bit quantization \citep{ashkboos2024quarot}; SpinQuant
learns rotations for improved weight/activation/KV quantization; DuQuant uses
outlier-aware rotations and transformations; and SmoothRot, IsoQuant,
PolarQuant, and ButterflyQuant explore related smoothing, isotropic, polar, or
structured-orthogonal transforms for low-bit LLMs
\citep{liu2024spinquant,lin2024duquant,huang2025smoothrot,kim2025isoquant,zhang2025polarquant,li2025butterflyquant}.
% Integration note: final bibliography verification for the rotation cluster
% is tracked in the citation audit.
ParoQuant is the closest current rotation baseline for our setting because it
targets reasoning LLM inference with pairwise rotations and W4A16-style
weight-only PTQ \citep{liang2025paroquant}.  Our related-work boundary is
therefore narrow: if a ParoQuant-style baseline recovers the audited gap, the
paper should become a diagnosis-and-baseline-contextualization result; if it
does not, the negative result is still scoped to the tested small/hybrid
long-decode regime rather than to rotation quantization in general.

\paragraph{Long-decode reasoning quantization.}
Reasoning models expose quantization failures that can be invisible in
short-horizon perplexity or chat benchmarks.  PMPD varies precision across
prefill, early decode, and later decode phases \citep{chen2025pmpd}, while
PM-KVQ targets long chain-of-thought KV-cache quantization with progressive
mixed precision and long-context calibration \citep{liu2025pmkvq}.  DecDEC is
the closest activation-channel systems prior: it dynamically selects salient
channels at each decode step and fetches residual compensation for those
channels \citep{park2025decdec}.  Our differentiation from DecDEC has four
axes.  First, \emph{horizon}: DecDEC studies short decode windows, while this
paper measures channel-set behavior at reasoning-scale decode positions.
Second, \emph{model class}: DecDEC targets Transformer LLMs, while our
intended surface includes small reasoning LLMs and hybrid reasoning models.
Third, \emph{architecture}: DecDEC's mechanism is a reactive Transformer
activation-channel residual path; OutlierMigrate measures whether static or
lightly position-conditioned protected sets remain stable in Transformer,
Mamba-hybrid, and parallel-hybrid layouts.  Fourth, \emph{benchmarks}: DecDEC
uses standard language/task evaluation for a low-bit systems method, whereas
our surface is long reasoning traces with preregistered drift and W4A16
recovery gates.  DecDEC therefore preempts any claim that changing salient
channels are new; our contribution must be framed as long-horizon measurement,
mechanism isolation, and audited failure or recovery of specific protection
policies.

\paragraph{Mamba and SSM quantization.}
State-space models add recurrence-specific quantization issues.  Quamba2
argues that Mamba activations exhibit channel-order preservation and
activation persistence, supporting offline per-channel quantization for
Mamba1/Mamba2 \citep{chiang2025quamba2}.  MambaQuant uses
variance-aligned rotations for Mamba-family PTQ \citep{xu2025mambaquant}.
QMamba reports highly dynamic hidden-state variation in vision SSMs and
introduces Temporal Group Quantization \citep{li2025qmamba}; OuroMamba
likewise identifies dynamic activation-outlier variation over time steps in
vision Mamba and uses lightweight dynamic outlier detection
\citep{ramachandran2025ouromamba}.  These works make the novelty boundary
strict: dynamic outliers are not new to Mamba broadly.  The open question for
this paper is whether the language-reasoning, long-decode, W4A16 protected-set
regime follows Quamba2-style persistence, vision-Mamba-style dynamic outlier
behavior, or a mixed failure mode that depends on architecture and horizon.

\paragraph{Recent adjacent adaptive quantizers.}
Several 2025--2026 methods reinforce the need to separate channel identity,
precision assignment, tensor type, and online adaptation.  ChanMix and MixKVQ
adapt precision across sensitive KV-cache channels for long-context or
reasoning workloads \citep{wen2026chanmix,zhang2025mixkvq}.  LAQuant uses
layer-wise lookahead loss for large reasoning model quantization
\citep{choi2026laquant}.  QEP explicitly propagates quantization error through
layer-wise PTQ to compensate accumulated error \citep{arai2025qep}.  FBQuant
uses feedback signals for LLM quantization, CMPQ allocates mixed precision in
a channel-wise pattern, and FlexQuant switches precision dynamically during
generation \citep{liu2025fbquant,chen2024cmpq,liu2025flexquant}.
% Integration note: final bibliography verification for adjacent adaptive
% quantizers is tracked in the citation audit.
These papers are adjacent rather than substitutable baselines.  KV-cache
methods act on keys and values; layer-wise error-propagation methods address
cross-layer reconstruction; and dynamic precision methods can change bitwidth
without necessarily changing which activation channels are protected.  Their
shared lesson is that static calibration is often too weak for modern
long-context inference.

\paragraph{Counter-results and terminology.}
SLQ is an important counterpoint: in a statistically near-lossless regime,
quantization error need not compound across token positions
\citep{helcig2026slq}.  We therefore do not assume that long decode alone
causes accumulating degradation.  The paper's mechanism split must keep three
possibilities separate: stale channel identity, insufficient protected budget,
and true long-horizon error accumulation.  Similarly, ``hot channels'' and
``TTQ'' require disambiguation.  Hot channels in recent NVFP4/pretraining
work refer to persistently extreme channels during training, not necessarily
the drifting W4A16 inference-time channels measured here.\footnote{We use
``hot channels'' only descriptively for high-magnitude channels on a measured
trace.  This is distinct from hot-channel pretraining terminology, where a
channel may be persistently extreme across training data or training steps.}
TTQ can mean classical Trained Ternary Quantization or recent test-time
quantization; this paper refers to the latter only when discussing
inference-time online calibration \citep{koikeakino2026ttq}.\footnote{TTQ is
ambiguous.  In this paper, ``test-time TTQ'' means activation-aware
test-time quantization for prompt-specific online calibration, not Trained
Ternary Quantization.}
% Integration note: TTQ is retained as a terminology note unless the final
% paper expands its test-time quantization comparison.
